% !PRINTED 592
\weberpage{592}
Sucht man den gr\"o\ss ten gemeinschaftlichen Teiler von $\boldsymbol{A}_\mu$ und
$\boldsymbol{A}_{\mu_1}$, so erh\"alt man eine ganze Funktion $\boldsymbol{A}_{\mathfrak{m}}$, deren Koeffizient % [source sic: the printed singular Koeffizient conflicts with the plural predicate]
ganze Zahlen in $\mathfrak{L}$ sind, deren Wurzeln die Gr\"o\ss en (9) sind, und
wenn man endlich $\boldsymbol{A}_{\mathfrak{m}}$ von Faktoren befreit, die es mit irgend
einem $\boldsymbol{A}_{\mathfrak{m}'}$ gemein hat, in dem $\mathfrak{m}'$ ein echter Teiler von $\mathfrak{m}$ ist, so
erh\"alt man eine ganze Funktion:
\weq{10}{T_{\mathfrak{m}}(x) = x^{\nu} - \tau_{\nu-1}x^{\nu-1} + \tau_{\nu-2}x^{\nu-2} - \cdots \pm \tau_0,}
deren Wurzeln nur die unter den Gr\"o\ss en (8) sind, in denen $\xi$
relativ prim zu $\mathfrak{m}$ ist.
% [source sic: the roots have just been identified as the quantities (9), but this sentence calls them quantities (8)]

Der Grad $\nu$ dieser Funktion ist nach Bd.\ II, \S\,168 zu bestimmen.
Wir wollen sie die \gesperrt{Idealteilungsfunktion} nennen. Ihre Koeffizienten
sind \gesperrt{ganze Zahlen} des K\"orpers $\mathfrak{L}$, und es ist
speziell
\weq{11}{\tau_0 = \overset{\xi}{\Pi}\,\boldsymbol{s}\!\left(\frac{2\alpha\xi}{\mu}\right),}
worin $\xi$ ein vollst\"andiges System inkongruenter, zu $\mathfrak{m}$ teilerfremder
Zahlen durchl\"auft.

Die Wurzeln von $T_{\mathfrak{m}}(x)$ sind gleich und entgegengesetzt, und
wir erhalten eine Gleichung von Graden $1/2\,\nu$ f\"ur die Gr\"o\ss e
\[
\boldsymbol{s}\!\left(\frac{2\alpha\xi}{\mu}\right)^{2}.
\]
Durch Adjunktion dieser Gr\"o\ss en zu dem K\"orper $\mathfrak{L}$ entsteht
ein K\"orper $\mathfrak{T}_{\mathfrak{m}}$, den wir gleichfalls Teilungsk\"orper nennen. Er
ergibt sich aus dem Teilungsk\"orper des \S\,154 durch Adjunktion
von $x$. Die Gruppe dieses K\"orpers ist in der Gruppe der Zahlen
$\xi\ (\mathrm{mod}\ \mathfrak{m})$ enthalten, wie leicht aus den Multiplikationsformeln
folgt.

Wir bemerken noch, da\ss{} sich der allgemeine Teilungsk\"orper
$\mathfrak{T}_{\mathfrak{m}}$ durch Anwendung des Additionstheorems auf den Fall
zur\"uckf\"uhren l\"a\ss t, wo $\mathfrak{m}$ eine Potenz eines Primideals in $\Omega$ ist.

\begin{center}
\textbf{\S\ 159. Primideale.}
\end{center}

6. \gesperrt{Ist $\eta$ eine gebrochene Zahl des K\"orpers $\Omega$, die in
reduzierter Form den ungeraden Idealnenner $\mathfrak{a}$ hat, und ist $\mathfrak{m}$
ein zu $\mathfrak{a}$ relativ primes Ideal, so ist die ganze Zahl $\boldsymbol{s}(\eta)$ relativ
prim zu $\mathfrak{m}$.}

Denn man nehme in $\Omega$ eine durch $\mathfrak{a}$ teilbare, zu $2\mathfrak{m}$ teilerfremde
 ganze Zahl $\mu$ an. Dann ist $\mu\eta$ eine ganze Zahl. In dem
Produkt \S\,157, (23)
%CONT>>
